\input{../tex-inputs/latex-leseansicht-vorspann}

\section[Arthur Schnitzler an Gustav Schwarzkopf, {[}7. 8. 1903?{]}]{L04200 Arthur Schnitzler an Gustav Schwarzkopf, {[}7. 8. 1903?{]}}
\nopagebreak\mylabel{L04200v}
\rehead{ }\normalsize\beginnumbering\briefempfaengerindex{Schwarzkopf, Gustav@\textsc{Schwarzkopf, Gustav}!zzzSchnitzler, Arthur@\emph{von Arthur Schnitzler}!1903-08-071@{{[}7. 8. 1903?{]}}|(be}
\toendnotes[C]{\smallbreak\pagebreak[2]}
\correspDesc{Versand  durch Arthur Schnitzler am [7. 8. 1903?] in Wien
\newline{}Erhalt  durch Gustav Schwarzkopf im Zeitraum [7. 8. 1903 – 10. 8. 1903?] in Wien}\toendnotes[C]{\smallbreak}
\Standort{CUL, Schnitzler, B 96.}
\physDesc{Karte, 200 Zeichen
\newline{}Handschrift: Bleistift, deutsche Kurrent}\toendnotes[C]{\smallbreak}
\pstart
           \noindent{}{\pb}lieber Guſtav, morgen Samſtag bei jedem Wetter ſind wir
                  Abend 8 in Hietzing\oindex{XIII., Hietzing@\textbf{XIII., Hietzing}, \emph{Verwaltungsgebiet}|pw}{ }\label{K_L04200-1v}\edtext{neues
               Reſtaurant \textsc{Kuffner\oindex{Wien@\textbf{Wien}!XIII., Hietzing@\textbf{XIII., Hietzing}!Ottakringer Bräu@\textbf{Ottakringer Bräu}, \emph{Bierhaus}|pw}}}{\lemma{\textnormal{\emph{neues
               Restaurant Kuffner}}}\Cendnote{\textnormal{Das erlaubt die Datierung, da Schnitzler in den Jahren 1903
                  und 1904 nur einmal an einem Samstag im Ottakringer Bräu (Kuffner)\oindex{Wien@\textbf{Wien}!XIII., Hietzing@\textbf{XIII., Hietzing}!Ottakringer Bräu@\textbf{Ottakringer Bräu}, \emph{Bierhaus}|pwk} in Hietzing\oindex{XIII., Hietzing@\textbf{XIII., Hietzing}, \emph{Verwaltungsgebiet}|pwk}
                  war, am 8. 8. 1903. Schwarzkopf\pwindex{Schwarzkopf, Gustav 7.\,11.\,1853 Wien – 13.\,11.\,1939 ebd.@\textsc{Schwarzkopf, Gustav} (7.\,11.\,1853 Wien – 13.\,11.\,1939 ebd.), \emph{Schriftsteller}|pwk} kam nicht.
               }}}\label{K_L04200-1} bei der Dampfhalteſtelle Neue Welt\oindex{Wien@\textbf{Wien}!XIII., Hietzing@\textbf{XIII., Hietzing}!Haltestelle Dommayergasse@\textbf{Haltestelle Dommayergasse}|pw}. Hugo\pwindex{Hofmannsthal, Hugo von 1.\,2.\,1874 Wien – 15.\,7.\,1929 Rodaun@\textsc{Hofmannsthal, Hugo von} (1.\,2.\,1874 Wien – 15.\,7.\,1929 Rodaun), \emph{Schriftsteller}|pw} läßt Sie bitten, we{\geminationn} Sie Zeit haben, zu kommen. Wir\pwindex{Schnitzler, Olga 17.\,1.\,1882 Wien – 13.\,1.\,1970 Lugano@\textsc{Schnitzler, Olga} (17.\,1.\,1882 Wien – 13.\,1.\,1970 Lugano), \emph{Schauspielerin, Sängerin}|pwv} nicht minder.\pend
           \pstart Ihr \spacefill\mbox{A.}\pend{}\selectlanguage{ngerman}\endnumbering\briefempfaengerindex{Schwarzkopf, Gustav@\textsc{Schwarzkopf, Gustav}!zzzSchnitzler, Arthur@\emph{von Arthur Schnitzler}!1903-08-071@{{[}7. 8. 1903?{]}}|)be}\mylabel{L04200h}
\begin{anhang}
\end{anhang}\newcommand{\dateiname}{L04200}\newcommand{\titel}{Arthur Schnitzler an Gustav Schwarzkopf, [7. 8. 1903?]}\newcommand{\editorInnen}{Herausgegeben von Jahnke, SelmaMüller, Martin Anton}\input{../tex-inputs/latex-leseansicht-abspann}
